\documentclass[11pt,reqno]{amsart}
\usepackage{iftex}
\ifPDFTeX
  \usepackage[T1]{fontenc}
  \usepackage[utf8]{inputenc}
  \usepackage{lmodern}
\else
  \usepackage[no-math]{fontspec}
\fi
\usepackage{amsmath}
\usepackage{amssymb}
\usepackage{amsfonts}
\usepackage{mathrsfs}
\usepackage{amsthm}
\usepackage{mathtools}
\usepackage{geometry}
\usepackage{enumitem}
\usepackage{longtable}
\usepackage{booktabs}
\usepackage{array}
\usepackage[dvipsnames]{xcolor}
\usepackage[pagebackref=true]{hyperref}
\usepackage{bookmark}
\usepackage{pdfcomment}

% Hover previews for internal references (no layout change).
\makeatletter
\let\CG@origref\ref
\let\CG@origeqref\eqref
\newcommand{\CG@gettip}[1]{\@ifundefined{CGTip@#1}{Reference: #1}{\csname CGTip@#1\endcsname}}
\DeclareRobustCommand{\CGref}[1]{\pdftooltip{\CG@origref{#1}}{\CG@gettip{#1}}}
\DeclareRobustCommand{\CGeqref}[1]{\pdftooltip{\CG@origeqref{#1}}{\CG@gettip{#1}}}
\expandafter\gdef\csname CGTip@sec:setting\endcsname{Section: Setting and notation}
\expandafter\gdef\csname CGTip@sec:lemma\endcsname{Section: The multiplicity lemma}
\expandafter\gdef\csname CGTip@sec:cert\endcsname{Section: Certificates from the defect}
\expandafter\gdef\csname CGTip@eq:defect\endcsname{Equation: u\_k = k - 1/2 - F(gamma\_k), ordinates counted with multiplicity}
\expandafter\gdef\csname CGTip@eq:Q\endcsname{Equation: Q\_\{2j\}(T,X) = (1/M\_T) int\_\{T\}\textasciicircum{}\{2T\} (N - F\_X)\textasciicircum{}\{2j\} dF\_X}
\expandafter\gdef\csname CGTip@thm:A\endcsname{Theorem - Exact defect identity (Theorem A of the clock manuscript): int (N-F)\textasciicircum{}\{2j\} F' dt = sum Psi\_j(u\_k) + boundary terms, Psi\_1(u) = u\textasciicircum{}2 + 1/12, Psi\_2(u) = u\textasciicircum{}4 + u\textasciicircum{}2/2 + 1/80}
\expandafter\gdef\csname CGTip@eq:Psi\endcsname{Equation: Psi\_1(u) = u\textasciicircum{}2 + 1/12, Psi\_2(u) = u\textasciicircum{}4 + u\textasciicircum{}2/2 + 1/80}
\expandafter\gdef\csname CGTip@lem:mult\endcsname{Lemma - Multiplicity lemma: if an ordinate has multiplicity m >= 2, its defects are u, u+1, ..., u+m-1 and sum (2(u+i))\textasciicircum{}\{2j\} >= m for every j >= 1; an ordinate of multiplicity 1 carries a simple zero on the critical line}
\expandafter\gdef\csname CGTip@thm:cert\endcsname{Theorem - Defect certificate: for every function F and every j >= 1, N\_0\textasciicircum{}s(T,2T) >= M\_T - sum\_\{k in I\_T\} (2u\_k)\textasciicircum{}\{2j\} - O(log T)}
\expandafter\gdef\csname CGTip@eq:cert\endcsname{Equation: N\_0\textasciicircum{}s(T,2T) >= M\_T - sum (2u\_k)\textasciicircum{}\{2j\} - O(log T)}
\expandafter\gdef\csname CGTip@cor:Q\endcsname{Corollary - Certificates in terms of Q: N\_0\textasciicircum{}s/M >= 4/3 - 4 Q\_2 - o(1) and N\_0\textasciicircum{}s/M >= 1 - 16 Q\_4 + 8 Q\_2 - 7/15 - o(1)}
\expandafter\gdef\csname CGTip@eq:cert2\endcsname{Equation: N\_0\textasciicircum{}s/M\_T >= 4/3 - 4 Q\_2 - o(1)}
\expandafter\gdef\csname CGTip@eq:cert4\endcsname{Equation: N\_0\textasciicircum{}s/M\_T >= 1 - 16 Q\_4 + 8 Q\_2 - 7/15 - o(1)}
\makeatother

\hypersetup{
  colorlinks=false,
  pdfborder={0 0 1},
  linkbordercolor={1 0 0},
  citebordercolor={0 0.3 1},
  urlbordercolor={0 0.6 0.6},
  filebordercolor={0 0.6 0.6},
  linktoc=all,
  bookmarksnumbered=true,
  bookmarksopen=true,
  bookmarksopenlevel=1,
  pdfstartview=FitH,
  pdftitle={Simple zeros on the critical line from the Euler-phase defect},
  pdfauthor={Rotkin Aynasac},
  pdfsubject={Riemann zeta function, simple zeros on the critical line, ordinal defect},
  pdfkeywords={Riemann zeta function, critical line, simple zeros, Euler product, pair correlation}
}
% Back-references in the bibliography.
\renewcommand*{\backref}[1]{}
\renewcommand*{\backrefalt}[4]{%
  \ifcase #1 %
    \relax
  \or
    \ (Cited on p.~#2.)%
  \else
    \ (Cited on pp.~#2.)%
  \fi}

% In \cite[note]{key} the note (e.g. a section number) is linked too.
\makeatletter
\let\CG@origcite\cite
\renewcommand*{\cite}[2][]{\def\CG@citekey{#2}%
  \ifx\relax#1\relax\CG@origcite{#2}\else\CG@origcite[#1]{#2}\fi}
\def\@cite#1#2{{\@citestyle[\citeform{#1}\if@tempswa, \hyper@link{cite}{cite.\CG@citekey}{#2}\fi]}}
\makeatother

\geometry{margin=1in}
\numberwithin{equation}{section}
\setcounter{tocdepth}{3}
\setlength{\emergencystretch}{1.5em}

\theoremstyle{plain}
\newtheorem{theorem}{Theorem}[section]
\newtheorem{proposition}[theorem]{Proposition}
\newtheorem{lemma}[theorem]{Lemma}
\newtheorem{corollary}[theorem]{Corollary}
\theoremstyle{definition}
\newtheorem{definition}[theorem]{Definition}
\newtheorem{hypothesis}[theorem]{Hypothesis}
\newtheorem{heuristic}[theorem]{Heuristic}
\theoremstyle{remark}
\newtheorem{remark}[theorem]{Remark}

\begin{document}
\title{Simple zeros on the critical line from the Euler-phase defect}
\author{Rotkin Aynasac}
\email{rotkivanyasac@gmail.com}
\date{20 October 2026}
\begin{abstract}
For any function $F$ and the ordinates $\gamma_k$ of the nontrivial zeros counted with multiplicity, put $u_k=k-\frac12-F(\gamma_k)$. An ordinate of multiplicity $m\ge2$ forces $u_{k+i}=u_k+i$ along its indices. From this we prove, unconditionally and for every $j\ge1$, that the number of simple zeros on the critical line in $(T,2T]$ is at least $M_T-\sum_k(2u_k)^{2j}-O(\log T)$. With the exact defect identity of the prime-corrected clock manuscript this becomes $N_0^s/M_T\ge\frac43-4Q_2-o(1)$, where $Q_2$ is the mean square of $N-F_X$ with respect to $dF_X$, $F_X=N_0+S_X$ and $S_X$ the Euler phase of the primes $p\le X$. A positive proportion of simple zeros on the critical line follows from $\limsup_TQ_2(T,T^c)<\frac13$ for one fixed $c$.
\end{abstract}
\maketitle
\tableofcontents

\section{Setting and notation}\label{sec:setting}
Let $\rho=\beta+i\gamma$ run over the nontrivial zeros of $\zeta(s)$, and let $\gamma_1\le\gamma_2\le\cdots$ be their ordinates in $(0,\infty)$, counted with multiplicity. Let $N(t)$ be the right-continuous counting function, $N_0(t)=\frac1\pi\Theta_0(t)+1$ with $\Theta_0(t)=\frac t2\log\frac t{2\pi e}-\frac\pi8$, and $\mathcal L=N_0'$. For $X\ge2$ let $S_X(t)=-\frac1\pi\sum_{p\le X}\arg(1-p^{-1/2-it})$ and $F_X=N_0+S_X$. For a real function $F$ put
\begin{equation}\label{eq:defect}
u_k=k-\tfrac12-F(\gamma_k),
\end{equation}
and, for $j\ge1$, with $\mathcal I_T=\{k:T<\gamma_k\le2T\}$ and $M_T=\#\mathcal I_T$,
\begin{equation}\label{eq:Q}
Q_{2j}(T,X)=\frac1{M_T}\int_T^{2T}\bigl(N(t)-F_X(t)\bigr)^{2j}\,dF_X(t).
\end{equation}
$N_0^s(T,2T)$ is the number of distinct zeros with $T<\gamma\le2T$, $\beta=\frac12$ and multiplicity $1$.

We use one result of the prime-corrected clock manuscript (version 11).

\begin{theorem}[Exact defect identity, {\cite[Thm.~A]{PCC11}}]\label{thm:A}
Let $F$ be $C^1$ on an open interval containing $[\gamma_{k_1},\gamma_{k_2}]$ and $u_k$ as in \CGeqref{eq:defect}. For every $j\ge1$,
$\int_{\gamma_{k_1}}^{\gamma_{k_2}}(N-F)^{2j}F'\,dt=\sum_{k_1<k<k_2}\Psi_j(u_k)+\frac{(\frac12+u_{k_1})^{2j+1}+(\frac12-u_{k_2})^{2j+1}}{2j+1}$, where
\begin{equation}\label{eq:Psi}
\Psi_1(u)=u^2+\tfrac1{12},\qquad \Psi_2(u)=u^4+\tfrac12u^2+\tfrac1{80}.
\end{equation}
\end{theorem}

\section{The multiplicity lemma}\label{sec:lemma}
\begin{lemma}\label{lem:mult}
Let $F$ be any real function and $u_k$ as in \CGeqref{eq:defect}.
\begin{enumerate}[label=\textup{(\alph*)}]
\item If an ordinate $\gamma$ has multiplicity $m\ge2$ in the list $(\gamma_k)$, with indices $k,\dots,k+m-1$, then $u_{k+i}=u_k+i$ for $0\le i<m$, and for every $j\ge1$
\[
\textstyle\sum_{i=0}^{m-1}\bigl(2(u_k+i)\bigr)^{2j}\ \ge\ m .
\]
\item If $\gamma$ has multiplicity $1$, the zero with ordinate $\gamma$ is simple and lies on the critical line.
\end{enumerate}
\end{lemma}

\begin{proof}
(a) $F(\gamma_{k+i})=F(\gamma)$ for all $i$, so $u_{k+i}-u_k=i$. Put $\phi(x)=(2x)^{2j}$, even, convex, $\phi(\pm\frac12)=1$. For $m=2$, $|u|+|u+1|\ge1$ and convexity give $\phi(u)+\phi(u+1)\ge2\phi(\frac12)=2$. For $m\ge3$, at most one of the $m$ numbers $u+i$ lies in $(-\frac12,\frac12)$, and the others have absolute values at least $\frac12,\frac12,\frac32,\frac32,\dots$ in increasing order; hence the sum is at least $2+2\cdot3^{2j}\lfloor(m-2)/2\rfloor+(m-2-2\lfloor(m-2)/2\rfloor)3^{2j}\ge m$.
(b) If $\rho=\beta+i\gamma$ with $\beta\ne\frac12$, then $1-\bar\rho=1-\beta+i\gamma$ is another zero with the same ordinate; a multiple zero repeats its ordinate. In both cases the multiplicity of $\gamma$ is at least $2$.
\end{proof}

\section{Certificates from the defect}\label{sec:cert}
\begin{theorem}\label{thm:cert}
For every real function $F$, every $j\ge1$ and $T\ge2$,
\begin{equation}\label{eq:cert}
N_0^s(T,2T)\ \ge\ M_T-\sum_{k\in\mathcal I_T}(2u_k)^{2j}-O(\log T),
\end{equation}
with an absolute implied constant.
\end{theorem}

\begin{proof}
By Lemma~\CGref{lem:mult}(b) the zeros of $\mathcal I_T$ that are not counted by $N_0^s$ lie on ordinates of multiplicity $m\ge2$; by part (a) each such group of $m$ indices satisfies $m\le\sum_{i}(2u_{k+i})^{2j}$. Summing over the groups contained in $\mathcal I_T$ and adding the nonnegative terms of the other indices gives \CGeqref{eq:cert}; the groups that meet the ends of the window contain $O(\log T)$ zeros because $N(t+1)-N(t)\ll\log t$.
\end{proof}

\begin{corollary}\label{cor:Q}
For $F=F_X$ with $2\le X\le T$, as $T\to\infty$,
\begin{align}
\frac{N_0^s(T,2T)}{M_T}&\ \ge\ \frac43-4\,Q_2(T,X)-o(1),\label{eq:cert2}\\
\frac{N_0^s(T,2T)}{M_T}&\ \ge\ 1-16\,Q_4(T,X)+8\,Q_2(T,X)-\frac7{15}-o(1).\label{eq:cert4}
\end{align}
\end{corollary}

\begin{proof}
Apply Theorem~\CGref{thm:A} on $\mathcal I_T$ (the boundary terms and the change from $[\gamma_{k_1},\gamma_{k_2}]$ to $[T,2T]$ are $O((\log T)^{2j+1})=o(M_T)$, because $|N-F_X|\ll\log T+\sqrt X$ is not needed: the integrand is bounded on the two end pieces by $(\log T+\sigma_X)^{2j}(\mathcal L+M_X)$ and their length is $O(1)$; for $X\le T$ this is $o(M_T)$ only when $\sqrt X\,(\log T)^{2j}=o(T\log T)$, which holds). With \CGeqref{eq:Psi}, $\langle u^2\rangle=Q_2-\frac1{12}+o(1)$ and $\langle u^4\rangle=Q_4-\frac12\langle u^2\rangle-\frac1{80}+o(1)=Q_4-\frac12Q_2+\frac7{240}+o(1)$. Insert into \CGeqref{eq:cert} with $j=1$ and $j=2$.
\end{proof}

\begin{remark}
By \cite[Thm.~B]{PCC11}, $Q_2(T,T^c)=\frac1{2\pi^2}\log\frac1c+O(1+\sqrt{\log(1/c)})$ for $2\le T^c\le T^{1/2}$ with an absolute constant that is not computed, so \CGeqref{eq:cert2} reduces a positive proportion of simple zeros on the critical line to the inequality $\limsup_T Q_2(T,T^c)<\frac13$ for one fixed $c$.
\end{remark}

\begin{thebibliography}{99}
\bibitem{PCC11}
The prime-corrected clock and the zeros of the Riemann zeta function, manuscript, version 11 (2026), repository LawCasanya, \href{https://github.com/vcasana1/LawCasanya}{github.com/vcasana1/LawCasanya}.

\end{thebibliography}
\end{document}
